
\subsubsection{3.1 Overview}

The methodology is motivated by the suppressor confounding hypothesis: high-FAIR datasets tend to be newer, and newer datasets have had less time to accumulate runs. Three core design choices follow: proxy FAIR scoring from OpenML quality metrics, propensity-matched pair construction, and survival analysis of TTFR.

\subsubsection{3.2 Data Collection}

We access the OpenML dataset corpus via the openml-python API (version 0.14.0), collecting 5,000 active tabular datasets with machine-computed quality metrics and run timestamps. A key API limitation was discovered during execution: the bulk \texttt{list_datasets()} endpoint does not return \texttt{upload_date}. Consequently, the planned post-2018 cohort filter could not be enforced; all active datasets are included. This limitation is discussed in Section 6.

\subsubsection{3.3 Proxy FAIR Score Construction}

The F-UJI tool [Devaraju and Huber, 2021] was unavailable in our pipeline environment. FAIR compliance is instead operationalized via a \textbf{proxy Findable score} derived from OpenML machine-computed quality metrics:

\begin{table}[htbp]
\centering
\begin{tabular}{lll}
\toprule
FAIR Sub-criterion & Proxy Operationalization & Weight \\
\midrule
F1: Persistent Identifier & Binary: dataset has a DOI or persistent URI & 0.25 \\
F2: Rich Metadata & Normalized metadata richness (proportion of non-null quality fields) & 0.50 \\
F3: Searchable & Binary: indexed in the OpenML search API & 0.25 \\
\bottomrule
\end{tabular}
\end{table}

Datasets with proxy score ≥ 0.5 are classified as high-FAIR (n = 720, 14.4\%); score < 0.5 as low-FAIR (n = 4,280, 85.6\%). This proxy captures only structural and computational characteristics, not semantic FAIR dimensions such as Interoperability or Reusability. Because the proxy is a noisy, attenuated operationalization of the intended IV, any significant finding under the proxy likely implies a stronger true effect under actual F-UJI sub-criteria scores.

\subsubsection{3.4 Propensity Score Matching}

Propensity scores are estimated using logistic regression on three covariates: creation year quartile (inferred from dataset ID ordering, as upload dates are unavailable), task type (classification/regression/other), and dataset size decile (log of instances × features). We apply 1:1 nearest-neighbor matching with caliper = 0.8 SD of the PS logit. A caliper of 0.8 SD is relaxed relative to the production target of 0.2 SD, and was necessary to obtain sufficient matched pairs from the smoke-test cohort of n = 200. Covariate balance is verified via standardized mean difference (SMD) < 0.1 for all covariates post-matching.

The necessity of matching is demonstrated empirically: unadjusted KM log-rank p = 0.583 (null result) versus matched KM p = 0.0053. This p-value reversal on the same datasets and same outcome variable is the central demonstration that suppressor confounding operates in this data.

\subsubsection{3.5 Survival Analysis of Time-to-First-Run}

\textbf{Primary outcome:} Time-to-first-run (TTFR) — days from dataset creation to first recorded experimental run on OpenML. TTFR captures discovery friction that FAIR Findable compliance theoretically reduces.

\textbf{Statistical methods:} Kaplan-Meier log-rank test (primary gate: p < 0.05, direction: high-FAIR < low-FAIR); Cox proportional hazards regression (secondary gate: HR > 1.2). The Schoenfeld residuals test flagged a potential proportional hazards (PH) violation in the smoke-test cohort. The reported Cox HR = 3.159 should therefore be interpreted as an average effect estimate; time-varying hazard models are warranted at production scale.

\subsubsection{3.6 Ablation and Sensitivity Design}

Three ablations assess sub-criteria specificity:
\begin{itemize}
\item \textbf{Ablation A:} Aggregate FAIR threshold (binary: overall quality score ≥ 0.5) as IV, replacing the Findable sub-criteria IV.
\item \textbf{Ablation B:} Accessible sub-criteria as IV.
\item \textbf{Ablation C:} Relaxed matching caliper (wider than primary setting).
\end{itemize}

Three sensitivity analyses test observation window robustness: 365 days (SA-2), 730 days (primary), and 1095 days (SA-3).

\subsubsection{3.7 Implementation Details}

Python 3.9; key libraries: lifelines 0.27.4 (Kaplan-Meier, Cox PH), scikit-learn 1.2.0 (propensity score estimation), openml 0.14.0, scipy 1.10.0. Random seed: 42. Unit tests: 30/30 passing for h-m1; 24/24 passing for h-m2. No GPU required.

